\section{Battery Monitoring}

This chapter contains information on the planned built-in battery monitoring and no-load voltage measuring.

\subsection{Introduction}

One of our goal for the project is to make it possible for the nodes and the gateway to be placed in an area for a long period of time without the need for maintenance. So we put a lot of emphasis on being as power efficient as possible to increase the amount of time a gateway or node can be up and running. To ensure that we could determine the battery life of the gateways and nodes, we wanted to implement a battery monitoring system in our project that would be able to transmit the current battery voltage values to the cloud as a raw voltage value.

\subsection{Battery Voltage Levels}

In the project we want to use a 3.7V LiPo battery as the primary power source. However, the voltage of 3.7V does not directly match the operating voltage requirements of our gateway and nodes, so a voltage conversion is required.

The gateway in our system requires a stable 5V supply to operate. To achieve this, a boost converter is used to bring the battery voltage up to 5V.

On the other hand, our nodes only need 3.3V to operate. A step-down converter is used to ensure that the battery voltage is converted to a stable 3.3V.

\subsection{Interfacing with the Microcontroller}

To monitor the system's battery voltage, we have implemented a function to read the raw battery voltage values using the analogue-to-digital converter pin P0.29. However, as the battery voltage is connected directly to the ADC pin of the microcontroller for the nodes, we needed to implement a voltage divider circuit to safely connect the battery to the ADC. This circuit scales the battery voltage level down to a level that the microcontroller can safely measure.

\subsection{No-load measurement}

Because direct measurement of the battery voltage while the system is operating can lead to inaccuracies on the voltage reading, we have integrated a no-load measurement function into the circuit to provide a more accurate reading of the battery voltage level.

To achieve this, we implemented a transistor and a capacitor in the circuit, then by sending a digital output of 5V to the transistor, we could cut off the connection to the microcontroller's battery for a short time and measure the battery level directly without load, while the microcontroller was still powered by the capacitor in the meantime.

\subsection{Difficulties in the implementation}

During the development of the battery monitoring system, we encountered voltage level problems that prevented us from including this feature in the final prototype. 

The use of a step-up and step-down converter prevented us from achieving the stable voltage levels required for the gateways and nodes in combination with other parts used in the circuit at the same time.

Another major problem was the transistor used in the no-load voltage measurement circuit. While the transistor was essential for the no-load measurement, it introduced a significant 2V voltage drop when active.

Using the humidity sensor and battery monitoring circuit in a single design also caused problems, as using two different voltage dividers for sensor and battery monitoring meant that the microcontroller only received 1.7V. This made it inoperable.

